\section{Garbled circuits}
Imagine Alice and Bob each have some secret values $a$ and $b$, and
would like to jointly compute some function $f$ over their respective
inputs. Furthermore, they would like to keep their secret values hidden
from each other: if Alice and Bob follow the protocol honestly, they
should both end up learning the correct value of $f(a,b)$, but Alice
should not learn anything about $b$ (other than what could be learned
by knowing both $a$ and $f(a,b)$), and likewise for Bob.

Here is our problem setting, slightly more formally:
\begin{goal}{}{}
\begin{itemize}
\item Alice knows a secret bitstring $a$ of length $m$ bits.
\item Bob knows a secret bitstring $b$ of length $n$ bits.
\item $f$ is a binary circuit, which takes in $m + n$ bits, and runs
  them through some $k$ gates. The outputs of some of the gates are
  the public outputs of the circuit. Without loss of generality, let us
  also suppose that each gate in $f$ accepts either $1$ or $2$
  input bits, and outputs a single output bit.
\item Alice and Bob would like to jointly compute $f(a,b)$ without
  revealing their secrets to each other.
\end{itemize}
\end{goal}


Yao's garbled circuits\footnote{Introduced by Andrew C.\ Yao in 1982: ``Protocols for secure computations''}
is one of the most well-known 2PC protocols.
The protocol is quite clever,
and optimized variants of the protocol are
being implemented and used today.\footurl{https://github.com/privacy-scaling-explorations/mpz/tree/dev/crates/mpz-garble}

\subsection{Outline of solution}
Our solution will contain two key components:

\begin{itemize}
\item Alice constructs a \alert{garbled circuit} that takes in the value
  $b$ (whatever it is) and spits out $f(a,b)$.
  A garbled circuit, roughly speaking, is an ``encrypted'' circuit that takes
  encrypted input and creates encrypted output.
\item \alert{Oblivious transfer} is a protocol where Alice has two
  messages, $m_{0}$ and $m_{1}$. Bob can get exactly one of them,
  $m_{i}$, without letting Alice know what $i$ is. In this context,
  Alice ends up sending Bob a password for his input in a way that Bob
  does not learn the passwords for any other inputs, and Alice does not
  find out which password she sent to Bob.
\end{itemize}

In slightly more detail:
\begin{enumerate}
\item
  Whatever the function $f$ is, we will assume that it takes $m + n$
  bits of input $a_{1},\ldots,a_{m}$ and $b_{1},\ldots,b_{n}$, and
  that it is computed by some sort of circuit made of AND, OR and NOT
  gates.
\item
  Alice's first task is to ``plug her own inputs into this circuit
  $f$.'' The result will be a new circuit (you might call it
  $f_{a}$) that has just $n$ input slots for Bob's $n$ bits
  $b_{1},\ldots,b_{n}$.
\item
  Now, Alice is going to ``garble'' the circuit $f_{a}$. Once it is
  garbled, Bob will not be able to see how it works. He will only be
  able to plug his own input $( b_{1},\ldots,b_{n})$ into
  the circuit.
\item
  To prevent Bob from plugging other inputs in as well, a garbled
  circuit will require a ``password'' for each input Bob wants to plug
  in -- a different password for every possible input. If Bob has the
  password for $( b_{1},\ldots,b_{n})$, he can learn
  $f_a( b_{1},\ldots,b_{n}) = f(a,b)$, but he will not
  learn anything else about how the circuit works.
\item
  Now, Alice has all the passwords for all the possible inputs, but how
  can she give Bob the password for
  $( b_{1},\ldots,b_{n})$? Alice does not want to let Bob
  have any other passwords -- and Bob is not willing to tell Alice which
  password he is asking for. This is where we will use ``oblivious
  transfer.''
\end{enumerate}

We now flesh out this outline, starting with garbled circuits.

\subsection{Garbled gates}
Our garbled circuits are going to be built out of \alert{garbled gates}.
A garbled gate is like a traditional gate (like AND, OR, NAND, NOR),
except its functionality is hidden.

What does that mean? Let us say the gate has two input bits, so there
are four possible inputs to the gate: $(0,0),(0,1),(1,0),(1,1)$.
For each of those four inputs $(x,y)$,
there is a secret password $P_{x,y}$.
The gate $G$ will only reveal its value $G(x,y)$ if you
give it the password $P_{x,y}$.

Here is a natural approach to make a garbled gate. Choose a
\emph{symmetric-key}\footnote{Symmetric-key encryption is probably what
  you think of when you think of plain-vanilla encryption: You use a
  secret key $K$ to encrypt a message $m$, and then you (or someone
  else) need the same secret key $K$ to decrypt it.} encryption scheme
$\Enc$ and publish the following table:

\begin{center}
\begin{tabular}{|c|c|}
\hline
$(0,0)$ & $\Enc_{P_{0,0}}( G(0,0))$ \\
\hline
$(0,1)$ & $\Enc_{P_{0,1}}( G(0,1))$ \\
\hline
$(1,0)$ & $\Enc_{P_{1,0}}( G(1,0) )$ \\
\hline
$(1,1)$ & $\Enc_{P_{1,1}}( G(1,1))$ \\
\hline
\end{tabular}
\end{center}

If you have the values $x$ and $y$, and you know the password
$P_{x,y}$, you just go to the $(x,y)$ row of the table, look up
\[\Enc_{P_{x,y}} ( G(x,y) ),\] decrypt it, and
learn $G(x,y)$.

But if you do not know the password $P_{x,y}$, assuming
$\Enc$ is a suitably secure encryption scheme, you will
not learn anything about the value $G_{x,y}$ from its encryption.

\subsection{Chaining garbled gates}

The next step is to combine a bunch of garbled gates into a circuit. We
will need to make some small changes to the protocol.

\begin{enumerate}
\item
  To chain garbled gates together, we need to modify the output of each
  gate: In addition to outputting the bit $z = G_{i}(x,y)$, the
  $i$-th gate $G_{i}$ will also output a password $P_{z}$ that Bob
  can use at the next step.

  Now Bob has one bit coming in for the left-hand input $x$, and it
  came with some password $P_{x}^{\text{left}}$ -- and then another
  bit coming in for $y$, with some password $P_{y}^{\text{right}}$.
  To get the combined password $P_{x,y}$, Bob concatenates the two
  passwords $P_{x}^{\text{left}}$ and $P_{y}^{\text{right}}$.
\item
  To keep the functionality of the circuit hidden, we do not want Bob to
  learn anything about the structure of the individual gates -- even the
  single bit he gets as output.

  This is an easy fix: Instead of having the gate output both the bit
  $z$ and the password $P_{z}$, we will have the gate just output
  $P_{z}$.

  \end{enumerate}

  But now how does Bob know what values to feed into the next gate? The
  left-hand column of the ``gate table'' needs to be indexed by the
  passwords $P_{x}^{\text{left}}$ and $P_{y}^{\text{right}}$, not by
  the bits $(x,y)$. But we do not want Bob to learn the other
  passwords from the table!

  Let us say this again. We want:

  \begin{itemize}
  \item
    If Bob knows both passwords $P_{x}^{\text{left}}$ and
    $P_{y}^{\text{right}}$, Bob can find the row of the table for the
    input $(x,y)$.
  \item
    If Bob does not know the passwords, he cannot learn them by looking
    at the table.
  \end{itemize}

  Of course, the solution is to use a hash function! So here is the new
  version of our garbled gate. For simplicity, we will assume it is an
  AND gate -- so the outputs will be (the passwords encoding) 0, 0, 0,
  1.


\begin{center}
\begin{tabular}{|c|c|}
\hline
  $\hash(P_{0}^{\text{left}},P_{0}^{\text{right}})$ &
  $\Enc_{P_{0}^{\text{left}},P_{0}^{\text{right}}}( P_{0}^{\text{out}})$ \\
  \hline
  $\hash(P_{0}^{\text{left}},P_{1}^{\text{right}})$ &
  $\Enc_{P_{0}^{\text{left}},P_{1}^{\text{right}}}( P_{0}^{\text{out}})$ \\
  \hline
  $\hash(P_{1}^{\text{left}},P_{0}^{\text{right}})$ &
  $\Enc_{P_{1}^{\text{left}},P_{0}^{\text{right}}}( P_{0}^{\text{out}})$ \\
  \hline
  $\hash(P_{1}^{\text{left}},P_{1}^{\text{right}})$ &
  $\Enc_{P_{1}^{\text{left}},P_{1}^{\text{right}}}( P_{1}^{\text{out}} )$ \\
  \hline
  \end{tabular}
\end{center}


\subsection{How Bob uses one gate}

Let us play through one round of Bob's gate-using protocol.

\begin{enumerate}
\item
  Suppose Bob's input bits are 0 (on the left) and 1 (on the right). Bob
  does not know he has 0 and 1 (but we do!). Bob knows his left password
  is some value $P_{0}^{\text{left}}$, and his right password is some
  other value $P_{1}^{\text{right}}$.
\item
  Bob takes the two passwords, concatenates them, and computes a hash.
  Now Bob has
  \[\hash(P_{0}^{\text{left}},P_{1}^{\text{right}}).\]
\item
  Bob finds the row of the table indexed by
  $\hash(P_{0}^{\text{left}},P_{1}^{\text{right}})$, and
  he uses it to look up
  \[\Enc_{P_{0}^{\text{left}},P_{1}^{\text{right}}}( P_{0}^{\text{out}}).\]
\item
  Bob uses the concatenation of the two passwords
  $P_{0}^{\text{left}},P_{1}^{\text{right}}$ to decrypt
  $P_{0}^{\text{out}}.$
\item
  Now Bob has the password for the bit 0 to feed into the next gate --
  but he does not know his bit is 0.
\end{enumerate}

So Bob is exactly where he started: he knows the password for his bit,
but he does not know his bit. So we can chain together as many of these
garbled gates as we like to make a full garbled circuit.

\subsection{How the circuit ends}

At the end of the computation, Bob needs to learn the final result. How?

Easy! The final output gates are different from the intermediate gates.
Instead of outputting a password, they will just output the resulting
bit in plain text.

\subsection{How the circuit starts}

This is trickier. At the beginning of the computation, Bob needs to
learn the passwords for all of his input bits. Let us frame the problem
for just a single bit.

\begin{itemize}
\item
  Alice has two passwords, $P_{0}$ and $P_{1}$.
\item
  Bob has a bit $b$, either $0$ or $1$.
\item
  Bob wants to learn one of the passwords, $P_{b}$, from Alice.
\item
  Bob does not want Alice to learn the value of $b$.
\item
  Alice does not want Bob to learn the other password.
\end{itemize}

This is where \emph{oblivious transfer} comes in, which we will see in
Section \ref{ot}.
